% =====================================================================
% research-review standard Beamer preamble  (XeLaTeX, 16:9, slate-blue)
% Copy this block verbatim at the top of papers/<slug>/slides/slide.tex,
% then add \begin{document} ... your frames ... \end{document}.
% Build: make slide PAPER=<slug>   (XeLaTeX 2-pass, out-of-tree → build/)
% =====================================================================
\documentclass[10pt, aspectratio=169]{beamer}

% --- Theme (clean, no external theme packages) -----------------------
\usetheme{default}
\usecolortheme{default}
\useinnertheme{circles}
\setbeamertemplate{navigation symbols}{}

% --- Fonts: XeLaTeX + fontspec, sans-serif body ----------------------
% fontspec defaults to Latin Modern (loaded by file from the texmf tree),
% so this needs NO fontconfig/system font registration — portable on any
% TeX Live / MacTeX. Do NOT \setmainfont{Latin Modern Roman}: the by-name
% lookup needs fontconfig and fails on stock MacTeX.
\usepackage{fontspec}
\usefonttheme{professionalfonts}      % keep math glyphs in CM, not sans
\renewcommand{\familydefault}{\sfdefault}   % Latin Modern Sans body
% For a CJK deck, add:  \usepackage{kotex}  (or set a CJK font via fontspec)

% --- Packages --------------------------------------------------------
\usepackage{xcolor}
\usepackage{graphicx}
\usepackage{amsmath}
\usepackage{amssymb}
\usepackage{mathrsfs}
\usepackage{booktabs}
\usepackage{array}
\usepackage{colortbl}                 % \rowcolor for table row highlight
\usepackage[font=scriptsize]{caption}
\usepackage[most]{tcolorbox}
\usepackage{tikz}
\usetikzlibrary{arrows.meta, positioning}   % for conceptual flow diagrams
\usepackage{pifont}                   % \ding{51}/\ding{55} check/cross marks for comparison tables

% --- Palette (3-5 colors, WCAG AA contrast) --------------------------
\definecolor{accent}{HTML}{1F3A5F}      % deep slate-blue — titles, rules, links
\definecolor{accentSoft}{HTML}{5C7AA0}  % muted slate — subtitles
\definecolor{ink}{HTML}{1A1A1A}         % body text
\definecolor{muted}{HTML}{6B6B6B}       % captions, metadata
\definecolor{rulec}{HTML}{D8DEE6}       % thin separator rules
\definecolor{boxbg}{HTML}{F4F6FA}       % tcolorbox / block background tint
\definecolor{crimson}{HTML}{780016}     % rare emphasis only (diagrams)
% --- Semantic keyword emphasis: positive = blue, negative = magenta --------
\definecolor{posblue}{HTML}{0066FF}     % positive-wording emphasis (blue)
\definecolor{negmag}{HTML}{D6009C}      % negative-wording emphasis (magenta)
\newcommand{\good}[1]{{\color{posblue}\bfseries #1}}  % wrap positive keyword
\newcommand{\warn}[1]{{\color{negmag}\bfseries #1}}   % wrap negative keyword
\newcommand{\blue}{\color{posblue}\bfseries}  % legacy switch -> positive (now bold)
\newcommand{\hl}{\color{posblue}\bfseries}    % legacy standout -> positive blue

% --- Beamer colors ---------------------------------------------------
\setbeamercolor{normal text}{fg=ink}
\setbeamercolor{structure}{fg=accent}
\setbeamercolor{frametitle}{fg=accent}
\setbeamercolor{framesubtitle}{fg=accentSoft}
\setbeamercolor{title}{fg=accent}
\setbeamercolor{subtitle}{fg=accentSoft}
\setbeamercolor{section in toc}{fg=accent}
\setbeamercolor{itemize item}{fg=accent}
\setbeamercolor{itemize subitem}{fg=accentSoft}
\setbeamercolor{enumerate item}{fg=accent}
\setbeamercolor{block title}{fg=white, bg=accent}
\setbeamercolor{block body}{fg=ink, bg=boxbg}
\setbeamercolor{caption name}{fg=accent}

\setbeamerfont{frametitle}{series=\bfseries, size=\large}
\setbeamerfont{framesubtitle}{series=\mdseries, size=\small}
\setbeamerfont{title}{series=\bfseries, size=\LARGE}
\AtBeginDocument{\footnotesize}   % single body baseline for the whole deck

% --- Footline: short-title (left) / page (right) ---------------------
\setbeamertemplate{footline}{%
  \leavevmode%
  \hbox{%
    \begin{beamercolorbox}[wd=\paperwidth, ht=2.4ex, dp=1.5ex, leftskip=1.2em, rightskip=1.2em, sep=0pt]{}%
      {\color{muted}\scriptsize\insertshorttitle}%
      \hfill%
      {\color{accent}\scriptsize$\bullet$}\hspace{0.6em}%
      {\color{muted}\scriptsize\insertframenumber\,/\,\inserttotalframenumber}%
    \end{beamercolorbox}%
  }%
  \vspace{2pt}%
}

% --- Itemize bullets -------------------------------------------------
\setbeamertemplate{itemize item}{\small$\blacktriangleright$}
\setbeamertemplate{itemize subitem}{\small$\triangleright$}
\setbeamertemplate{itemize subsubitem}{\scriptsize\textendash}

% --- Section divider: one full-bleed page per section ----------------
\AtBeginSection[]{%
  \begin{frame}[plain, noframenumbering]
    \vfill
    \begin{center}
      {\color{muted}\scriptsize\MakeUppercase{Section \insertsectionnumber}}\\[8pt]
      {\color{accent}\Large\bfseries\insertsectionhead}\\[12pt]
      {\color{accent}\rule{0.22\paperwidth}{0.8pt}}
    \end{center}
    \vfill
  \end{frame}%
}

% --- tcolorbox defaults: inline title (box width = content width) ----
\tcbset{
  enhanced,
  colframe=accent!55,
  colback=boxbg,
  coltitle=accent,
  boxrule=0.4pt,
  arc=2pt,
  left=5pt, right=5pt, top=3pt, bottom=3pt,
  fonttitle=\bfseries\small,
  title style={left color=accent!8, right color=accent!8},
  separator sign={\,:},
}

\hypersetup{colorlinks=true, linkcolor=accent, citecolor=accent, urlcolor=accent}

\DeclareMathOperator*{\argmax}{argmax}
\DeclareMathOperator*{\argmin}{argmin}
